\section*{9. Vorläufige Gliederung}

% keine Abkürzungen bzw. vorher definieren
%   o	z.B. im Überblick als Kapitel
%   o	später dann in der Gliederung-Story darauf eingehen

\begin{enumerate}
    % 1.
    \item Einleitung
    % 2.
    \item Analyse des Technikstandes in Bezug auf das Testing sowie die Optimierung durch Modernisierung von graphischen Oberflächen an Lega\-cy-Systemen
    \begin{enumerate}
        \item Aufbau und Eigenschaften von Testsystemen mit Fokus auf gewachsene Legacy-Umgebungen
        \item Ansätze zur Modernisierung von Legacy-Systemen
        \begin{enumerate}
            \item Methoden zur Modernisierung anhand unterschiedlicher Eingriffstiefe in das Legacy-System
            \item Methoden und Strategien zur Migration von Legacy-Systemen
            \item Technische Schulden und Strategien zur Tilgung
        \end{enumerate}
        \item Grundlagen zur Gestaltung nutzerfreundlicher und moderner Oberflächen (UI/UX) in Bezug zur Erstellung von Dashboards
        \begin{enumerate}
            \item DIN EN ISO 9241-110 zur Gebrauchstauglichkeit und Ergonomie interaktiver Systeme
            \item Standardisierter Fragebogen nach ISO 9241-110
            \item moderne Tools und Empfehlungen zur Erstellung von Dashboards
        \end{enumerate}
    \end{enumerate}
    % 3.
    \item Dashboards als potenzielle Lösung von Effizienz- und Übersichtlichkeitsproblemen in Legacy-Test-Systemen
    \begin{enumerate}
        \item Herausforderungen in der Erstellung von Dashboards und der Arbeit mit Legacy- und Code-Testing-Systemen
        \item Fokussierung auf die Problemstellung der Effizienz und Übersichtlichkeit verteilter Testsysteme
        \item Beitrag moderner Dashboards zur Effizienzsteigerung und verbesserten Übersichtlichkeit von Legacy-Testsystemen
        \item Methodisches Vorgehen zur Entwicklung eines Dashboards als moderne Anzeige der Testpoolergebnisse
        \begin{enumerate}
            \item Analyse des Ist-Zustandes der Testpool-Ergebnisanzeige zur Identifikation bestehender Anforderungen
            \item Entwicklung eines Dashboards als moderne Ergebnisübersicht
            \item Usability Test zur Validierung der gestellten Anforderungen an das Dashboard
            \item Nutzerbefragung zur Validierung der Benutzererfahrung im Rahmen des Usability-Tests
        \end{enumerate}
    \end{enumerate}
    % 4.
    \item Entwicklung eines Dashboards als moderne Anzeige der Testpoolergebnisse anhand zuvor analysierter bestehender System- und Nutzeranforderungen
    \begin{enumerate}
        \item Anforderungsanalyse anhand des bestehenden Legacy-Systems und den Nutzerbedürfnissen
        \item Evaluation und Auswahl des Dashboard-Tools
        \item Überblick über den strukturellen Aufbau und Inhalt des Dashboards
        \item technische Umsetzung des Datenflusses zur Dashboardanzeige
    \end{enumerate}
    % 5. 
    \item Validierung des entwickelten Dashboards in Bezug auf die Effizienz, Nutzererfahrung und -freundlichkeit im Vergleich zum Legacy-System
    \begin{enumerate}
        \item Durchführung der Validierungsverfahren
        \item Umsetzung eines Usability-Tests zur Überprüfung der Nutzerfreundlichkeit und Effizienz des entwickelten Dashboards im Vergleich zum Legacy-System
        \begin{enumerate}
            \item Erstellung eines Usability-Tests zur Überprüfung der identifizierten Anforderungen wie Funktionsvollständigkeit, Nutzerfreundlichkeit und Effizienz
            \item Zusammenfassung der Testergebnisse im Rahmen des Legacy-Systems
            \item Zusammenfassung der Testergebnisse im Rahmen des Dashboards
        \end{enumerate}
        \item Umsetzung einer Nutzerbefragung als zusätzliche Validierung der Nutzerfreundlichkeit und Gewährleistung der Nutzerzufriedenheit im Vergleich zum Legacy-System
        \begin{enumerate}
            \item Erstellung eines Fragebogens nach ISO 9241-110 zur Validierung der gegebenen Anforderungen wie User-Experience und Nutzerzufriedenheit
            \item Zusammenfassung der Ergebnisse des Fragebogens in Bezug auf das Legacy-System
            \item Zusammenfassung der Ergebnisse des Fragebogens in Bezug auf das Dashboard
        \end{enumerate}
    \end{enumerate}
    % 6.
    \item Auswertung der Ergebnisse und Methodenkritik zur Bewertung der Eignung, Effizienz und User-Experience der Dashboard-Lösung im Vergleich zum Legacy-System
    \begin{enumerate}
        \item Auswertung der Ergebnisse des Usability Tests in Bezug auf die funktionelle Eignung und Effizienz des Dashboards im Vergleich zum Legacy-System
        \begin{enumerate}
            \item Analyse der Erfüllung der Funktionsanforderungen des entwickelten Dashboards
            \item Bewertung der Effizienz des Dashboards durch Gegenüberstellung der Testergebnisse
        \end{enumerate}
        \item Auswertung der Ergebnisse der Nutzerbefragung zur Validierung der Testergebnisse sowie der Bewertung der User-Experience des Dashboards
        \begin{enumerate}
            \item Gegenüberstellung der Befragungsergebnisse des Legacy-Systems und Dashboards zur Validierung der vorangegangenen Dashboardbewertung
            \item Auswertung der Nutzerbefragung zur Bewertung der User-Experience des Dashboards
        \end{enumerate}
        \item Kritische Auseinandersetzung mit der verwendeten Methodik zur Modernisierung von Legacy-Systemen
        \item Kritische Auseinandersetzung mit der verwendeten Methodik zur Validierung der gewählten Dashboard-Lösung
        \item Abschließende Bewertung der Dashboard-Lösung zur Modernisierung eines Legacy-Systems
    \end{enumerate}
    % 7.
    \item Schlussfolgerungen, Limitationen und Ausblick
    \begin{enumerate}
        \item Zusammenfassung der wichtigsten Ergebnisse und Erkenntnisse
        \item Grenzen des entwickelten Dashboards und mögliche Weiterentwicklungen
        \item Ausblick auf weitere Möglichkeiten zur Modernisierung von Legacy-Systemen
    \end{enumerate}
\end{enumerate}
